% !TEX root = ../monograph.tex
\chapter{Decomposition into phases}
\label{ch15:phases}

A translation-invariant correlation family can be a mixture of states
with different long-distance behavior. We separate this mixture using
operators on the zero-energy subspace of its reflection Hilbert space.
A phase will be a character of a commutative operator algebra, not a
set of gauge configurations. Its correlations are obtained by evaluating
compressed whole words.

The inputs are the common field spaces $F_d$ and separated correlation
family of Chapter~\ref{ch13:continuum}, bounded slab operators with
their adjoint and concatenation laws, and orthogonal covariance of that
same family. The latter is stronger than the original lattice
hypercubic covariance. We also require the bounded real lattice
approximations and corrected moment estimates for separating families specified below.
The assertions retain these analytic hypotheses; a decomposition theorem
for an abstract state alone does not establish them.

\section{The zero-energy subspace}
\label{ch15:ground}

Use time coordinate four. Write $\mathcal H^{\rm OS}$ for the mixed
reflection Hilbert space, $\Omega$ for its unit vacuum,
$E(t)=e^{-tH}$ for the heat semigroup, and $U(y)$ for spatial
translations. Let $P_0$ be the projection onto
$\mathcal H_0=\ker H$. The spectral theorem gives
$E(t)\to P_0$ strongly as $t\to\infty$.

Every spatial translation acts trivially on $\mathcal H_0$.
Here the orthogonal covariance input is essential. For two separated
clusters put $F_{v,w}(x)=S(v\,\tau_xw)$. At sufficiently large time
separation this is a matrix element of $U(\vec x)E(x^4-c)$ between
fixed vectors. Hence a fixed temporal displacement changes it by a
quantity tending to zero, uniformly in arbitrary spatial displacements.
If $b^4>0$, choose an orthogonal map $K$ with $Kb=|b|e_4$.
Then $(Ke_4)^4=b^4/|b|>0$. Rotate the clusters and apply the temporal
statement along $tKe_4$. This proves
$F_{v,w}(te_4+b)-F_{v,w}(te_4)\to0$.
For spatial $y$ use $b=y+e_4$ and remove the extra unit time shift.
Taking strong heat limits and using the density of word vectors gives
\begin{equation}
                              U(y)P_0=P_0.
 \label{ch15:spatial-ground}
\end{equation}
This conclusion concerns the ground space of the mixture; it does not
say that this space is one dimensional.

\section{Compression and separation of clusters}
\label{ch15:compression}

For a word $w$ placed inside a slab of width $h$, let $B_h(w)$ be its
bounded insertion operator and define
\[
 M(w)=P_0B_h(w)|_{\mathcal H_0}.
\]
Increasing the width appends a heat operator, and shifting the word in
time inserts one on the other side. Both disappear after compression.
Spatial placement disappears by \eqref{ch15:spatial-ground}. Thus
$M(w)$ is defined for any distinct-point word independently of its
slab placement, with
\[
 M(\varnothing)=1,\qquad M(w)\Omega=P_0[w],\qquad
                  S(w)=\langle\Omega,M(w)\Omega\rangle.
\]
For positive-time words also put $Y_w=B_h(w)|_{\mathcal H_0}$,
which maps into the full reflection Hilbert space. The slab adjoint
and concatenation laws give
\begin{equation}
 M(\theta v\,w)=Y_v^*Y_w,
 \qquad M(\theta v)=M(v)^*.
 \label{ch15:whole-gram}
\end{equation}

If successive time gaps between translated words tend to infinity,
their correlation tends to the corresponding ordered product of
compressed operators:
\begin{equation}
 \lim_{j\to\infty}S(\tau_{z_1(j)}w_1\cdots\tau_{z_k(j)}w_k)
       =\langle\Omega,M(w_1)\cdots M(w_k)\Omega\rangle,
 \quad (z_{i+1}(j)-z_i(j))^4\to\infty.
 \label{ch15:cluster}
\end{equation}
The spatial parts of the displacements can be arbitrary. Indeed,
if $\psi_j\to\psi$ and $t_j\to\infty$, then
$U(y_j)E(t_j)\psi_j\to P_0\psi$ by contraction, strong heat
convergence and \eqref{ch15:spatial-ground}. Iterating this fact through
the bounded slab operators proves \eqref{ch15:cluster}.

\begin{proposition}[The compression algebra]
\label{ch15:algebra}
The norm-closed unital star algebra $\mathfrak A$ generated by all
$M(w)$ on $\mathcal H_0$ is commutative and separable. The vector
$\Omega$ is cyclic and separating for it.
\end{proposition}
\begin{proof}
We give the particular geometric commutativity argument. Place four
words at
\[
 0,\quad j(e_4+2e_1),\quad j(2e_4+e_1),\quad3j(e_4+e_1).
\]
Let $s$ exchange coordinates one and four. This reverses the time order
of the two middle clusters, as shown in Figure~\ref{ch15:fig-cluster-exchange}.
Covariance and \eqref{ch15:cluster} therefore exchange
the two middle operators, initially with that orthogonal map applied
to all four words. Apply the same cluster limit to four words placed
on the ray $j(e_4+e_1)$, which the exchange fixes. This removes the
common transformation of the labels and internal positions. Thus
\[
 \langle\Omega,M(w_1)[M(w_2),M(w_3)]M(w_4)\Omega\rangle=0.
\]
\begin{figure}[!htbp]
\centering
\begin{tikzpicture}[x=0.87cm,y=0.87cm,font=\small,
  axis/.style={->,>=Stealth,black!65,thin},
  centre/.style={circle,fill=ymInk,inner sep=0pt,minimum size=4.3pt},
  word/.style={fill=white,inner sep=1pt,outer sep=2pt}]
  \foreach \shift/\heading in {0/{Exchange of the middle clusters},6.5/{Centres on the fixed ray}}{
    \begin{scope}[xshift=\shift cm]
      \foreach \k in {1,2,3}{
        \draw[ymGrid,thin] (0,\k) -- (3.25,\k);
        \draw[ymGrid,thin] (\k,0) -- (\k,3.25);
      }
      \draw[axis] (-0.18,0) -- (3.8,0) node[right,black] {$x^1/j$};
      \draw[axis] (0,-0.18) -- (0,3.72) node[above,black] {$x^4/j$};
      \draw[ymInk,dashed] (0,0) -- (3.35,3.35);
      \node[centre] at (0,0) {}; \node[word,below right] at (0,0) {$w_1$};
      \node[centre] at (3,3) {}; \node[word,below right] at (3,3) {$w_4$};
      \node at (1.75,-0.92) {\heading};
    \end{scope}
  }
  \draw[<->,>=Stealth,ymAccent,thick,shorten <=3.5pt,shorten >=3.5pt]
    (2,1) -- (1,2);
  \node[centre] at (2,1) {}; \node[word,right] at (2,1) {$w_2$};
  \node[centre] at (1,2) {}; \node[word,left] at (1,2) {$w_3$};
  \begin{scope}[xshift=6.5cm]
    \node[centre] at (1,1) {}; \node[word,below right] at (1,1) {$w_3$};
    \node[centre] at (2,2) {}; \node[word,below right] at (2,2) {$w_2$};
  \end{scope}
\end{tikzpicture}
\caption[Exchanging two clusters in the phase algebra]{Each point gives the displacement of a whole word. Read the
  points from bottom to top for the original time order and from left to
  right for the time order after the coordinate exchange $s$. On the left
  the order changes from $w_1,w_2,w_3,w_4$ to $sw_1,sw_3,sw_2,sw_4$.
  On the right $s$ fixes all four centres. Here the cluster limit gives
  the same vacuum expectation with $s$ removed from each word in the
  reordered product.}
\label{ch15:fig-cluster-exchange}
\end{figure}

The vectors $M(v)\Omega=P_0[v]$ span a dense subspace of $\mathcal H_0$;
reflected words give the adjoints in the left slot. The bounded
commutator is therefore zero. Density proves cyclicity. If
$A\Omega=0$ for $A\in\mathfrak A$, commutativity gives $AB\Omega=0$
for all $B\in\mathfrak A$, hence $A=0$. Norm continuity of slab
operators, finite-dimensional label spaces, and countable dense
position sets give a countable generating family, hence separability.
\end{proof}

Compression is not multiplication of finite-separation words. In fact
\begin{equation}
 M(\theta v\,w)-M(v)^*M(w)=Y_v^*(1-P_0)Y_w.
 \label{ch15:compression-defect}
\end{equation}
The right side usually survives at finite separation. It is the
whole-word identity \eqref{ch15:whole-gram}, rather than a product of
scalar phase expectations, that will give phase reflection positivity.

\FloatBarrier
\section{Characters and the representing measure}
\label{ch15:characters}

Let $\Xi$ be the space of unital star characters of $\mathfrak A$.
The commutative Gelfand representation identifies
$\mathfrak A\simeq C(\Xi)$ isometrically; see the original
\cite[Section~2, Lemma~1]{GelfandNaimark1943}. The compact space $\Xi$
is metrizable because $\mathfrak A$ is separable. The positive normalized
vacuum functional has a unique representing probability measure
$\nu$ by the Riesz theorem; for this application see
\cite[Theorems~2.4.1 and~2.8.2]{Landsman1998}.
It has full support. Otherwise a nonzero continuous function supported
in a nonempty open null set would correspond to an $A\ne0$ with
$\|A\Omega\|^2=\int|\xi(A)|^2\,d\nu=0$, contradicting separation.

Define
\begin{equation}
 S^\xi(w)=\xi(M(w)),\qquad
 S(w)=\int_\Xi S^\xi(w)\,d\nu(\xi).
 \label{ch15:mixture}
\end{equation}
Every fixed separated value is continuous in $\xi$. This is a
representation of the mixed correlations by characters of the ground
algebra. Vacuum uniqueness in these phases will be proved separately.

We shall repeatedly use the following consequence of full support.
If $A\in\mathfrak A$ and
$\langle\Omega,(A^*A)^j\Omega\rangle\le C^{2j}$ for all $j$, then
\begin{equation}
                        \|A\|=\sup_{\xi\in\Xi}|\xi(A)|\le C.
 \label{ch15:moment-norm}
\end{equation}
This is the usual limit of $L^{2j}$ norms on a probability space,
combined with full support for a continuous function. In applications,
the moment bound comes by taking widely separated copies of the same
lattice expression, applying the chessboard bound at fixed copy number,
and then taking the ordered limits and the cluster limit.
The order of these operations matters.

\section{Symmetries in each phase}
\label{ch15:phase-symmetry}

Symmetry of a mixture need not imply symmetry of every component.
Here an additional lattice variance estimate gives the required
operator identity. Write $\vartheta$ for linear time reflection on
positions and labels, without conjugation. For a generic real grid word $w$ (distinct projections in each
coordinate),
set $Z=M(w)-M(\vartheta w)$. The adjoint law gives $Z^*=-Z$, so
\[
 a:=\langle\Omega,Z^2\Omega\rangle=-\|Z\Omega\|^2\le0.
\]
Corollary~\ref{ch13:bounded-word}, using the complete residual
estimate of Lemma~\ref{ch13:coarse-residual-supplier}, gives the
following bounded real approximation. For widely translated lattice copies $f_i$ of
$w-\vartheta w$, and every $\varepsilon>0$, there are real $g_i$
with $|g_i|\le R_\varepsilon$, uniformly in cutoff, volume and copy
position, such that
\[
 |\langle g_ig_j\rangle-\operatorname{Re}\langle f_if_j\rangle|
                                      \le\varepsilon\quad(i<j)
\]
on the ordered cutoff tails at fixed translation and fixed pair.
There is no corresponding assumed bound on $f_i^2$.
Variance positivity for $q+1$ copies gives
\[
 0\le\left\langle\left(\sum_{i=0}^qg_i\right)^2\right\rangle
 \le(q+1)R_\varepsilon^2+
              2\sum_{i<j}(\operatorname{Re}\langle f_if_j\rangle+\varepsilon).
\]
First take the ordered cutoff limits at fixed $q$ and translations,
then send their separation to infinity. Equation~\eqref{ch15:cluster}
gives $0\le(q+1)[R_\varepsilon^2+q(a+\varepsilon)]$.
Send $q\to\infty$ and then $\varepsilon\downarrow0$.
Thus $a=0$, and separation of $\Omega$ gives $Z=0$.
The cofinal grids are chosen after the degree and lower-level request.
Generic configurations are dense in the full separated configuration
space. Norm continuity and multilinearity
extend the identity to every word.

Orthogonal covariance of the mixture induces simultaneous invariance
of compressed product moments whenever a cluster-separating ray has
positive time before and after rotation. Such a ray exists unless
time is reversed; compose with $\vartheta$ in that case. On the dense
vectors $M(w)\Omega$ this gives a unitary orthogonal action, under
which time reflection now acts as the identity. Its conjugates, all
hyperplane reflections, therefore act trivially as well. Since these
generate $O(4)$,
\begin{equation}
                     M(Kw)=M(w)\qquad(K\in O(4)).
 \label{ch15:phase-invariant}
\end{equation}
Character evaluation proves symmetry in every phase. The proved lattice
approximation estimate is indispensable to this conclusion.

\section{Positivity, field identities and null spaces}
\label{ch15:fields}

For positive words $w_i$, \eqref{ch15:whole-gram} gives the positive
operator
\[
 \sum_{i,j}\overline{c_i}c_jM(\theta w_i\,w_j)
                   =\left(\sum_i c_iY_{w_i}\right)^*
                    \left(\sum_j c_jY_{w_j}\right)\ge0.
\]
It belongs to $\mathfrak A$. Positivity agrees with ambient operator
positivity in a unital C$^*$-subalgebra, and characters are positive;
see \cite[Section~2.6]{Landsman1998}. Hence each phase is reflection
positive, with a Gram form that need not have rank one.
Norm-smooth slab dependence gives smooth phase correlations and their
uniform separated bounds. Field inclusion, derivative and stress
identities pass to phases when written as operator identities, tested
against all exterior words before character evaluation. For a stress
shell the weights avoid collisions, so the norm-continuous operator
integral is defined. The expectation identity alone would not suffice.

A mixed null field has zero slab operator, by testing against the dense
word vectors. It is therefore null in every phase:
$N_d(S)\subset N_d(S^\xi)$. The reverse inclusion uses the quantitative separating families from the
same lattice family. Apply Proposition~\ref{ch11:n-mark-comparison}
with two marks per comparison cell. At a fixed sufficiently deep
comparison level $k$, its corrected coefficients for one and two insertions
have chessboard seminorm at most $B_k=H_k^{-1/32}$, uniformly in
cutoff and period. Choose a simultaneous coefficient limit for the
finitely many bounded lifts of the separating family at this fixed $k$, on a
refinement of the already selected correlation limits.
For every fixed integer $j$, put $2j$ copies of a corrected word
and its adjoint at widely separated translations. The chessboard
bound is $B_k^{2j}$: the two-mark estimate is used with arbitrary
replica census, not rerun with arity $2j$. Its constant is chosen
before $j$; only the cutoff and translation thresholds depend on
that fixed number. The repeated-word identity and the cluster
limit \eqref{ch15:cluster} identify the limiting product with
$\langle\Omega,(\Delta^*\Delta)^j\Omega\rangle$.
Taking successive deeper levels $k=k_\ell$ therefore supplies
compressed errors for one and two insertions $\Delta$ with
\begin{equation}
 \langle\Omega,(\Delta^*\Delta)^j\Omega\rangle
                  \le\gamma_\ell^{2j}\quad(j\ge1),
                  \qquad \gamma_\ell\longrightarrow0.
 \label{ch15:detector-moments}
\end{equation}
Here one may take $\gamma_\ell=B_{k_\ell}$, enlarging it by
a fixed finite row constant if necessary. The free row moments are
bounded, so $\gamma_\ell\to0$. These are estimates uniform in the number of reflected copies,
stronger than a bound on the mixed mean.
The finite reference rows use fixed free labels and convergent
noncollision configurations; their one-point means are uniformly bounded.
To make the pair correction explicit, if the reference one-point means are $s,s'$ and the reference connected pair is $c$, use
\[
 \Delta_1=M(u)-s1,\quad \Delta'_1=M(v)-s'1,\quad
 \Delta_2=M(uv)-sM(v)-s'M(u)+(ss'-c)1.
\]
Equation~\eqref{ch15:moment-norm} bounds each error by
$\gamma_\ell$ in every character. Therefore the phase pair differs
from $ss'+c$ by at most $\gamma_\ell(1+|s|+|s'|)$.
The phase matrices of correlation pairings converge uniformly in $\xi$ to the
injective free matrix of correlation pairings. Since they factor through
$V_d/N_d(S^\xi)$, its dimension is at least $\dim\mathcal M_d^{\rm Max}$.
The inclusion of null spaces gives the opposite bound through
Proposition~\ref{ch13:rank}. Consequently
\begin{equation}
                         N_d(S^\xi)=N_d(S)
                      \quad\hbox{for every }\xi.
 \label{ch15:same-fields}
\end{equation}
All phases have the same field filtration. Neither abstract character
decomposition nor positivity alone proves this equality.

\section{Selecting a phase with a unique ground vector}
\label{ch15:selection}

Construct the reflection Hilbert space and heat semigroup of each
phase using its separated hierarchy. Let $P_{0,\xi}$ denote its ground
projection. For a positive word $w$, the nonnegative heat matrix
elements decrease to $\|P_{0,\xi}[w]_\xi\|^2$ and are dominated by
the continuous zero-time Gram value. Dominated convergence and
\eqref{ch15:mixture} yield
\[
 \int\|P_{0,\xi}[w]_\xi\|^2\,d\nu=\|P_0[w]\|^2
                         =\int|S^\xi(w)|^2\,d\nu.
\]
Their difference is
$\|P_{0,\xi}[w]_\xi-S^\xi(w)\Omega_\xi\|^2$. It vanishes almost
everywhere. Rational positive configurations with basis labels form
one countable family dense in every phase reflection space, by Gram
continuity. Intersect its full-measure sets. Outside the resulting
null set,
\begin{equation}
                       \ker H_\xi=\CC\Omega_\xi.
 \label{ch15:ae-ground}
\end{equation}
Choose and retain one $\xi_*$ in this full-measure set. At this stage
we have not eliminated exceptional phases or shown that $\nu$ is a
point mass.

The slab bounds are uniform in all phases, including exceptional
ones. For the formal iterates $v_j=(A_{w^\dagger}A_w)^jv$, the mixed
maps satisfy
$Y_{v_j}=(B_h(w)^*B_h(w))^jY_v$. Thus
\[
 \|[v_j]_\xi\|^2=\xi(Y_{v_j}^*Y_{v_j})
                   \le\|Y_v\|^2\|B_h(w)\|^{4j}.
\]
The repeated-power argument gives bounded phase slab operators with
$\|B_h^\xi(w)\|\le\|B_h(w)\|$, and their identities and derivatives
follow on the dense word span. These uniform bounds are the input for
relativistic reconstruction and for the later removal of exceptional
phases. No spectral gap has been obtained here.

% !TEX root = ../monograph.tex
\section{Sharp separated bounds in every phase}
\label{ch15:sharp-phase-section}

Recall that \eqref{ch13:sharp-separated} bounds the undecomposed fixed-arity
correlation by the sum of degrees plus any positive loss in the exponent.
The estimate is also uniform in the phase $\xi$. To see this without
bounding phase evaluations by the undecomposed expectation of the same word,
apply the preceding chessboard argument to arbitrarily many separated
copies of that \emph{concatenated} word and their adjoints. The cell scales are
chosen from the classes modulo the tiling and the separations of one fixed word, not
from the total number of copies. Copies placed far apart occupy distinct
cells and have the same one-word factor on the right of the chessboard
inequality. Passing to the separated-cluster limit yields the same
bound to the corresponding power for the vacuum expectations of the powers of its compressed
operator $M(w)^*M(w)$. The norm formula of
\eqref{ch15:moment-norm} and the character bound then give
\begin{equation}
 \sup_{\xi\in\Xi}|S_n^\xi(\boldsymbol A;x)|
      \le C_{\boldsymbol A,n,\epsilon}r^{-\sum_i d_i-\epsilon}.
 \label{ch13:sharp-phase}
\end{equation}
Thus the fixed-arity scaling exponent is the sum of degrees plus an
arbitrarily small loss. No uniform constant as arity tends to infinity
is asserted; the repeated-word estimate remains the appropriate statement
for that different purpose.

